% ------------------------------------------------------------------
%  Capitolo 15 — Il corpo che impara
%  Parte IV
%  Ricerca di riferimento: ricerca 01 §11; 03 §3.5, §3.7
%  Voci bibliografiche principali: jenkins1924, rasch2013, yoo2007, hillman2008, dewar2012, smith2016
%  Epigrafe: ✔ verificata sul testo (vedi ricerca 03 e registro, ricerca 00).
%  Confermata come unica sede di questa epigrafe il 25/09/2026: il cap. 6
%  (due forze) usava lo stesso passo di Quintiliano ed è stato spostato su
%  un'altra fonte (Aristotele, radici amare/frutto dolce) per evitare il
%  doppione; qui il passo resta perché parla letteralmente di "una notte
%  interposta" (sonno). Il cap. 9 lo cita già in nota, rimandando qui.
%  Scheda del capitolo: ricerca/06_indice-ragionato.md, capitolo 15
% ------------------------------------------------------------------
\chapter{Il corpo che impara}
\label{cap:corpo}

\epigrafe[È straordinario a dirsi, e la ragione non è evidente, quanta solidità aggiunga una notte interposta.]{Mirum dictu est, nec in promptu ratio, quantum nox interposita adferat firmitatis.}{Quintiliano, Institutio oratoria XI, 2, 43}

\iniziale{Q}{uesto capitolo} \segnaposto{apertura: da scrivere — vedi la scheda nell'indice ragionato}.

\section{Il sonno che consolida}

\segnaposto{da scrivere}

\section{Il movimento}

\segnaposto{da scrivere}

\section{Stress e ansia}

\segnaposto{da scrivere}

\section{Pause vere}

\segnaposto{da scrivere}

\section{Cibi e integratori: poche prove}

\segnaposto{da scrivere}

\begin{daricordare}
\segnaposto{il principio del capitolo in due righe}
\end{daricordare}

\begin{domande}
  \item \segnaposto{domanda di recupero su questo capitolo}
  \item \segnaposto{domanda di applicazione a una materia che studi}
  \item \segnaposto{domanda di ripresa su un capitolo precedente}
\end{domande}

\begin{sintesi}
  \item \segnaposto{punto di sintesi}
\end{sintesi}

\finecapitolo
